\documentclass[
  blockchain,
  computerscience,
  chemistry,
  biochemistry,
  biology
]{genesisl1-paper}

\GenesisTitle{A Cross-Field Template for Reproducible Community Science}
\GenesisSubtitle{Equations, chemical compounds, code, wide figures and exact artifact references}
\GenesisRunningTitle{Cross-field reproducible community science}
\GenesisVersion{Version 3.0}
\GenesisStatus{COMMUNITY REVIEW}
\GenesisDate{August 2026}
\GenesisAuthors{GenesisL1 Community Publication Templates}
\GenesisAffiliations{Blockchain, computer science, chemistry, biochemistry and biology}
\GenesisCorrespondence{community@example.org}
\GenesisIdentifier{GL1-CPT-PAPER-3.0}
\GenesisRepository{https://github.com/JoMfN/genesisl1-community-publication-templates}
\GenesisFields{Cross-field profile}
\GenesisKeywords{community science; reproducibility; chemistry; software; provenance}

\GenesisDataIdentifier{doi:10.xxxx/example-dataset}
\GenesisCodeIdentifier{git:0123456789abcdef}
\GenesisModelIdentifier{model:v1.2.0}
\GenesisExecutionIdentifier{run:example-2026-08-01}
\GenesisContentHash{sha256:placeholder}
\GenesisIPFSCID{bafy...placeholder}
\GenesisIPFSGateway{Optional resolver supplied by the author}

\begin{document}

\begin{GenesisFrontMatter}
\begin{GenesisStructuredAbstract}
\AbstractBackground{
Reproducible community science requires a manuscript format that can express
mathematical models, chemical identities, biological nomenclature, software, and
verifiable artifacts without allowing infrastructure language to overwhelm the science.
}
\AbstractMethods{
A two-column \LaTeX{} class was implemented with structured abstracts, numbered
equations and compounds, language-aware code boxes, wide floats, numerical citations,
and optional content-addressed metadata.
}
\AbstractResults{
The example compiled with standard pdf\LaTeX{} and produced stable references to
\cref{eq:rate}, \cref{sch:compounds}, \cref{lst:python}, and \cref{fig:wide}.
}
\AbstractConclusion{
The resulting format supports cross-field scientific writing while keeping interpretation,
limitations, provenance, and future infrastructure distinct.
}
\end{GenesisStructuredAbstract}
\end{GenesisFrontMatter}

\section{Introduction}
Scientific manuscripts frequently connect mathematical analysis, molecular structures,
biological observations, software, and computational provenance. The template treats
these as compatible components rather than as reasons to impose one rhetorical model
on every author.

\section{Equations}
Numbered equations use standard \LaTeX{} semantics. For example, an elementary rate
model can be written as
\begin{equation}
  v = \frac{V_{\max}[S]}{K_{\mathrm m} + [S]},
  \label{eq:rate}
\end{equation}
and cited as \cref{eq:rate}. The equation remains human-readable in source and
professionally typeset in the document.

\section{Chemistry and biochemistry}
Chemical formulae use \texttt{mhchem}, for example
\begin{equation}
  \ce{ATP + H2O -> ADP + P_i + H+}.
  \label{eq:atp}
\end{equation}

\begin{scheme*}[t]
  \centering
  \begin{tabular}{c@{\hspace{18mm}}c}
    \chemfig{*6(-=-(-OH)=-=)} &
    \chemfig{*6(-=-(-NH_2)=-=)}\\[1mm]
    \compound{phenol-example} &
    \compound{aniline-example}
  \end{tabular}
  \caption{Example structures with persistent compound numbering. Compounds
  \compoundref{phenol-example} and \compoundref{aniline-example} can be cited
  consistently throughout the manuscript.}
  \label{sch:compounds}
\end{scheme*}

Biological names may be written as \taxon{Escherichia coli}, genes as \gene{lacZ},
and proteins as \protein{LacZ}. These lightweight commands can be adapted to the
nomenclature rules of a specific discipline.

\section{Computer science and blockchain}
Code is typeset with the standard \texttt{listings} package, so the default build does
not require shell escape.

\begin{GenesisCode}[
  language=Python,
  caption={A minimal deterministic transformation},
  label={lst:python}
]
from hashlib import sha256

def commit(payload: bytes) -> str:
    # Return a lowercase SHA-256 commitment.
    return sha256(payload).hexdigest()
\end{GenesisCode}

A Solidity fragment can be cited independently:

\begin{GenesisCode}[
  language=Solidity,
  caption={A content commitment stored by a contract},
  label={lst:solidity}
]
pragma solidity ^0.8.24;

contract ArtifactCommitment {
    mapping(bytes32 => address) public author;

    function register(bytes32 digest) external {
        require(author[digest] == address(0), "already registered");
        author[digest] = msg.sender;
    }
}
\end{GenesisCode}

The result in \cref{lst:python} describes a byte commitment. It does not establish the
quality, legality, safety, or scientific interpretation of the committed object.

\section{Wide figures}
Single-column figures use the ordinary \texttt{figure} environment. Important diagrams
may use \texttt{figure*} or the helper environment below to span both columns.

\begin{GenesisWideFigure}[t]
  \fbox{%
    \begin{minipage}{0.92\textwidth}
      \centering
      \vspace{7mm}
      \sffamily\bfseries\color{GenesisNavy}
      Scientific question
      \(\longrightarrow\)
      exact inputs
      \(\longrightarrow\)
      method or model
      \(\longrightarrow\)
      result
      \(\longrightarrow\)
      interpretation\\[2mm]
      \normalfont\small\color{GenesisInk}
      Each arrow denotes an explicit dependency, not automatic scientific validation.
      \vspace{7mm}
    \end{minipage}
  }
  \caption{A full-width figure spanning the two-column scientific body.}
  \label{fig:wide}
\end{GenesisWideFigure}

\section{Results}
The compiled example preserved the order and labels of equations, schemes, code
listings, and wide figures. All references were generated by \texttt{cleveref}.

\begin{GenesisEvidence}
The build produced a two-column paper in which the official GenesisL1 lockup appeared
only in the first-page front matter. Subsequent page headers contained text only.
\end{GenesisEvidence}

\section{Discussion}
Our results provide evidence that a single paper class can support the selected fields
without making every manuscript about blockchain infrastructure. This observation
points towards a broader community publication system, but suitability for a specific
journal or reporting guideline must still be evaluated separately.

\begin{GenesisInterpretation}
The main value is not visual uniformity by itself. It is the consistent distinction among
direct observations, interpretations, limitations, and exact supporting artifacts.
\end{GenesisInterpretation}

\begin{GenesisLimitations}
This example is a template demonstration rather than an empirical validation study.
Chemical structures generated from SMILES are intentionally left to a separate,
reviewable preprocessing step. The default manuscript build consumes committed
graphics or \texttt{chemfig} source and does not execute external converters.
\end{GenesisLimitations}

\section{Conclusion}
A supportive manuscript container can provide technical precision while leaving
scientific judgement, disciplinary conventions, and authorial voice intact.

\section*{Artifact record}
\PrintGenesisArtifactRecord

\GenesisEndMatterNotice

\begin{thebibliography}{2}
\bibitem{whitepaper}
GenesisL1.
\textit{A Sovereign Layer-1 for Verifiable Scientific Infrastructure}.
Technical Whitepaper, Version 1.0, 2026.

\bibitem{community}
GenesisL1 Community Publication Templates.
\textit{Scientific Tone and Authoring Guide}.
Version 3.0, 2026.
\end{thebibliography}

\end{document}
